\chapter{Anpassung und Lösungsansätze}\label{ch:anpassung}


\section{Vorschläge zur Anpassung der Unternehmensprozesse}\label{sec:anpassung-unternehmensprozesse}

Die in Kapitel~\ref{sec:gaps-auswirkungen} quantifizierten Gaps wurden zwar simulativ erzeugt, entsprechen aber Konfigurations- und Organisationsmustern, wie sie in Praxisberichten zu \gls{sap}-Anwender\-unternehmen wiederkehrend beschrieben werden \parencites[S.~3 ff.]{reinkemeyer2020processmining}[S.~21 ff.]{peters2019processmining}. Dieses Kapitel leitet ab, wie ein Unternehmen die gemessenen Ineffizienzen beheben kann, gegliedert nach organisatorischen Maßnahmen, Customizing und Erweiterungen.

Auf organisatorischer Ebene setzen drei Maßnahmen an den Wurzeln der Muster an. Erstens eine Stammdaten-Governance. Für Kundenstammdaten werden Pflichtattribute, Verantwortlichkeiten und ein Bereinigungsprojekt definiert, das Vorgehen orientiert sich an den bewährten Schritten der Datenmigration. Datenquellen identifizieren, klassifizieren, bereinigen und die Qualität kontinuierlich überwachen \parencite[Folie~324 ff.]{coletta2026_0206}. Vollständige Versanddaten und Incoterms entziehen den Liefersperren (G3) die Grundlage. Zweitens die Entlastung der Auftragserfassung, etwa durch elektronischen Datenaustausch mit den umsatzstärksten Kunden. Das reduziert Erfassungsfehler und damit die Änderungsschleife (G2). Drittens ein definierter Reklamationsprozess mit Rollen und Fristen, der die Gutschrifts- und Nachlieferungsfälle (G6) aus der informellen E-Mail-Bearbeitung in einen messbaren Ablauf überführt. Dass Stammdatenmängel der zentrale Engpasstreiber im \gls{o2c}-Prozess sind, bestätigt die Fallstudie von Adolfsson/Saleh, in der fehlerhafte Kundenstammdaten bei einem Kunden in allen Fällen manuelle Lieferterminänderungen und in rund einem Fünftel nachträgliche Preiskorrekturen erzwangen \parencite[S.~55 f.\ und 61]{adolfsson2026o2c}. Automatisierung, Datenintegration und fortgeschrittene Analytik gelten in der Literatur übergreifend als zentrale Hebel zur Verbesserung des \gls{o2c}-Prozesses in \gls{sap} \gls{sd} \parencite[S.~110 ff.]{nadukuru2024o2c}. Diese organisatorischen Maßnahmen sind Voraussetzung dafür, dass die anschließenden Systemeinstellungen überhaupt greifen. Eine automatisierte Kreditprüfung ist nur so gut wie die zugrunde liegenden Kreditstammdaten, und eine Pflichtfeldsteuerung für Versanddaten verhindert Liefersperren nur dann, wenn zugleich definiert ist, wer die fehlenden Attribute in welcher Frist pflegt. Organisatorische und technische Maßnahmen bilden also keine Alternative, sondern greifen ineinander. Praktisch beginnt eine Stammdaten-Governance mit der Benennung von Verantwortlichen je Datenobjekt, der Definition verbindlicher Pflichtfelder und Pflegestandards sowie einem einmaligen Bereinigungslauf, der Dubletten und veraltete Einträge entfernt; erst danach greifen die Pflichtfeldprüfungen im System, ohne den laufenden Betrieb zu blockieren.

\section{Möglichkeiten des Customizings im \gls{sap}-System}\label{sec:customizing}

Der Großteil der Gaps lässt sich ohne Programmierung im Customizing schließen, Tabelle~\ref{tab:massnahmenkatalog} ordnet jedem Gap die Maßnahme und den \gls{sap}-Bezug zu. Kern ist die Automatisierung der Kreditprüfung im \gls{sap} Credit Management. Bestandskunden mit gutem Zahlungsverhalten erhalten regelbasierte Scores und Limits, sodass Sperren nur noch bei tatsächlichen Risikofällen entstehen (G1). Die Unvollständigkeitsprüfung und Pflichtfeldsteuerung verhindert Aufträge mit fehlenden Versanddaten (G3), die Umstellung des Fakturavorrats auf tägliche Läufe beseitigt die künstliche Wartezeit vor der Rechnungsstellung (G4), und ein restriktives Berechtigungskonzept für manuelle Preisänderungen erzwingt die Pflege im Konditionsstamm (G5). Nachträgliche Auftragsänderungen (G2) werden durch eine entlastete, fehlerärmere Auftragserfassung gesenkt, insbesondere durch die \gls{edi}-Anbindung der umsatzstärksten Kunden (M6). Zu den typischen Customizing-Aktivitäten zählen dabei das Abbilden von Geschäftsprozessen, Grundeinstellungen und Berechtigungen \parencite[Folie~349]{coletta2026_0206}, Workflows für Freigaben (hier für Gutschriften) sind Teil der Systemkonfiguration \parencite[Folie~244]{coletta2026_0505}.

\begin{table}[!htb]
\centering
\caption{Maßnahmenkatalog: Gaps, Lösungsansätze und \gls{sap}-Bezug}\label{tab:massnahmenkatalog}
\small
\begin{tabularx}{\textwidth}{@{}>{\RaggedRight\arraybackslash}p{3.0cm} >{\RaggedRight\arraybackslash}p{2.8cm} >{\RaggedRight\arraybackslash}p{4.4cm} >{\RaggedRight\arraybackslash}p{3.9cm}@{}}
\toprule
\textbf{Maßnahme} & \textbf{Typ} & \textbf{\gls{sap}-Bezug} & \textbf{Adressiertes Gap /\allowbreak{} erwarteter Effekt} \\
\midrule
M1 Automatisierte Kreditprüfung & Customizing & \gls{sap} Credit Management: Scoring-Regeln, Limitprofile, Ausnahmeliste & G1: Sperrquote von 34,1 \% deutlich reduziert; Median-Liegezeit 1,9 AT entfällt für Bestandskunden \\
M2 Stammdatenbereinigung u. Pflichtfelder & organisatorisch + Customizing & Unvollständigkeitsprüfung, Pflichtfeldsteuerung, Geschäftspartnerpflege & G3: Liefersperrenquote von 12,3 \% Richtung Restrisiko; 3,9 $\rightarrow$ 0,8 AT bis Lieferbeleg \\
M3 Tägliche Fakturierungsläufe & Customizing & Fakturavorrat (VF04/\allowbreak{}VF06) als täglicher Hintergrundjob & G4: Liegezeit WA $\rightarrow$ Faktura von 2,6 auf ca. 0,5 AT \\
M4 Konditionspflege u. Berechtigungen & Customizing + organisatorisch & Preisfindung/\allowbreak{}Konditionstechnik, Rollenkonzept & G5: Fakturasperren (7,9 \%) weitgehend vermieden \\
M5 Workflow Reklamation/\allowbreak{}Gutschrift & Customizing/\allowbreak{}Workflow & Freigabeworkflow für Gutschriftsanforderungen & G6: definierte Durchlaufzeit, Transparenz \\
M6 \gls{edi}-Anbindung Top-Kunden & Erweiterung & \gls{idoc}-Eingang für Bestellungen (ORDERS) & G2: Änderungsquote (22,7 \%) sinkt; weniger Erfassungsfehler \\
\bottomrule
\end{tabularx}
\end{table}

Am Beispiel der Kreditprüfung (M1) lässt sich der Hebel konkretisieren. Im Ist-Zustand wird jeder Auftrag, der ein statisches Kreditlimit überschreitet, gesperrt und manuell freigegeben, unabhängig vom tatsächlichen Risiko. Das \gls{sap} Credit Management ordnet Kunden stattdessen Risikoklassen mit eigenen Prüfregeln und Limits zu, sodass Bestandskunden mit tadellosem Zahlungsverhalten automatisch freigegeben werden, während echte Risikofälle weiterhin geprüft werden. Die Sperre wird damit vom Regelfall zur Ausnahme, ohne den Schutz vor Forderungsausfällen aufzugeben. Zu beachten ist, dass mehrere Maßnahmen ineinandergreifen: M1 setzt gepflegte Kreditstammdaten voraus und der Verzicht auf manuelle Preise (M4) einen vollständigen Konditionsstamm. Die Stammdaten-Governance aus Kapitel~\ref{sec:anpassung-unternehmensprozesse} ist damit nicht eine Maßnahme unter mehreren, sondern das Fundament, auf dem die technischen Anpassungen aufbauen.

Über die Kreditprüfung hinaus greifen die übrigen Customizing-Maßnahmen an klar umrissenen Konfigurationspunkten an. Die Unvollständigkeitsprüfung wird so eingestellt, dass ein Auftrag ohne gepflegte Versanddaten gar nicht erst zur Lieferung freigegeben wird, was die Liefersperren an der Wurzel verhindert. Der Fakturavorrat wird von einem wöchentlichen auf einen täglichen Hintergrundjob umgestellt, und die Berechtigung, Preise manuell im Auftrag zu überschreiben, wird auf wenige Rollen beschränkt, sodass Sonderkonditionen künftig im Konditionsstamm gepflegt werden. Keine dieser Einstellungen erfordert eine Programmierung, sie nutzen ausschließlich vorgesehene Parameter des Standards, was Wartbarkeit und Release-Fähigkeit erhält.

\section{Modifikationen, Erweiterungen und kontinuierliches Monitoring}\label{sec:modifikationen-monitoring}

Eigenentwicklungen oder Modifikationen des \gls{sap}-Standards sind für keines der sechs Gaps erforderlich, im Einklang mit dem Leitsatz „So viel Standard wie möglich, so wenig Customizing wie nötig“ \parencite[Folie~248]{coletta2026_0505}. Als Erweiterung im Sinne des Enhancement-Typs des Process Mining empfiehlt sich jedoch ein kontinuierliches Prozessmonitoring. Die Extraktionslogik aus Kapitel~\ref{sec:datenstrukturen} wird als wiederkehrender Ladeprozess eingerichtet, sodass Kennzahlen wie Sperrquoten, Liegezeiten und Happy-Path-Anteil laufend überwacht und Maßnahmenwirkungen gemessen werden können. Process Mining wandelt sich damit von einer einmaligen Diagnose zu einem dauerhaften Steuerungsinstrument \parencite[S.~3 ff.]{reinkemeyer2020processmining}. Die verwendete Werkzeugkette aus Python und \gls{pm4py} zeigt, dass dafür nicht zwingend eine kommerzielle Plattform erforderlich ist, wie sie \gls{sap} mit dem Beispielinhalt „Process Mining on \gls{sap} \gls{s/4hana}“ in \gls{papm} für dasselbe Szenario anbietet \parencite[Abschnitt~2.1, S.~9]{sap2022samplecontent}.

Der Aufbau eines solchen Monitorings ist selbst ein kleines Datenprojekt. Die Extraktion muss regelmäßig und möglichst inkrementell erfolgen, die Kennzahlen sind in einem Dashboard zu verdichten, und für kritische Größen wie eine wieder steigende Sperrquote sind Schwellenwerte für Warnungen festzulegen. Damit wird aus der einmaligen Fit/Gap-Analyse ein geschlossener Regelkreis, in dem Abweichungen früh erkannt und gegengesteuert werden, bevor sie sich in Durchlaufzeit und Forderungslaufzeit niederschlagen. Konkret bieten sich als Kennzahlen die bereits in der Analyse verwendeten Größen an: die Sperrquoten je Sperrentyp, die Median-Liegezeiten je Prozessübergang, der Anteil konformer Fälle und die Durchlaufzeit bis zur Faktura. Werden sie zusätzlich je Verkaufsorganisation oder Kundengruppe aufgeschlüsselt, lassen sich Problembereiche gezielt eingrenzen, statt nur einen Gesamtdurchschnitt zu betrachten.

